\section{Заключение}
\label{sec:conclusion}

В настоящей работе предложен и реализован воспроизводимый подход к оценке
русскоязычных RAG-систем, основанный на переносе методологии RAGBench/TRACe на русский
язык. Работа объединяет несколько связанных результатов: построение русскоязычного
бенчмарка, обучение семейства токен-уровневых моделей-оценщиков, анализ их поведения на
разных конфигурациях и создание прикладного инструмента для применения обученных моделей.
В совокупности эти результаты формируют основу для более систематической оценки
русскоязычных RAG-систем как в исследовательских, так и в прикладных сценариях.

\subsection{Полученные результаты}

\begin{enumerate}
  \item \textbf{Построен русскоязычный бенчмарк.}
  Опубликован корпус RAGBench-RU\footnote{\url{https://huggingface.co/datasets/CMCenjoyer/ragbench-ru}},
  представляющий собой перевод 7 поднаборов исходного RAGBench на Qwen2.5-72B.
  При построении корпуса реализован конвейер, включающий структурные проверки
  корректности перевода, сохранности ключей и согласованности разметки.

  \item \textbf{Воспроизведена схема LLM-оценщика без дообучения.}
  Сопоставление с Qwen2.5-72B использовалось как опорный ориентир,
   а не как прямой эксперимент на одном корпусе.
  \item \textbf{Обучено семейство моделей-оценщиков.}
  Обучено около 80 конфигураций моделей-оценщиков на разных базовых архитектурах, длинах
  контекста, схемах весов функции потерь и вариантах выходного классификационного блока. Лучшая
  конфигурация - \textbf{ModernBERT-base, длина~2048, Simple~3:3:1} -
  достигает $\bar{F}_1 = 0.858$ на объединенном тестовом корпусе.

  \item \textbf{Проведен сравнительный экспериментальный анализ.}
  Показано, что на исследованной задаче:
    \begin{itemize}
      \item оптимальная длина контекста для лучших конфигураций находится вблизи 2048
      токенов, тогда как слишком короткий контекст приводит к потере информации, а
      чрезмерно длинный - к размыванию сигнала;
      \item различие между схемами весов 1:1:1 и 3:3:1 оказывается небольшим и не дает
      устойчивого практического преимущества;
      \item простая классификационная архитектура систематически устойчивее сложной,
      особенно на длинных контекстах;
      \item на построенном корпусе русскоязычные конфигурации показали более
      высокие значения $F_1$ по relevance и utilization, однако интерпретация
      этого различия ограничена влиянием перевода, токенизации и перенесённой
      разметки и не может быть сведена к чистому эффекту языка.
      \item дополнительная проверка на SberQuAD-QA показала, что обученные модели 
      демонстрируют переносимость
      по метке релевантности на внешнем русскоязычном
      наборе данных без дообучения и без подбора порогов.

    \end{itemize}

  \item \textbf{Классификатор как практическая альтернатива LLM-оценщику.}
  LLM-оценщик без дообучения на базе Qwen2.5-72B оценивался на исходном
  англоязычном корпусе RAGBench на уровне предложений (совпадение по
  предложениям, среднее число пропусков) и демонстрирует заметные расхождения
  с эталонной разметкой, включая склонность к избыточной разметке релевантных
  предложений. Обученный классификатор, в свою очередь, достигает высокого
  значения F1 на уровне токенов относительно эталонной разметки RAGBench-RU.
  Прямое сравнение Qwen и классификатора по единой метрике на уровне токенов
  и на одном корпусе в данной работе не проводилось: LLM-оценщик и классификатор
  запускались на разных корпусах (англоязычный RAGBench и русскоязычный
  RAGBench-RU соответственно) и в разных шкалах агрегации, поэтому сопоставление
  трактуется как сравнение с опорным ориентиром. Вместе с тем полученные
  результаты согласуются с ключевой идеей оригинальной статьи RAGBench о
  практической применимости компактного дообученного классификатора как
  существенно более дешёвой альтернативы LLM-оценщику.

  \item \textbf{Реализован прикладной сервис.}
  Разработан Streamlit-сервис, поддерживающий одиночный и пакетный режимы применения
  обученных моделей, настройку порогов и визуализацию токен-уровневой разметки. Сервис
  делает результаты работы доступными для прикладного использования, хотя не претендует
  на завершенную систему производственного развертывания.
\end{enumerate}

\subsection{Границы применимости и дальнейшее развитие}

Полученные результаты следует интерпретировать с учётом того, что RAGBench-RU
построен машинным переводом исходного англоязычного корпуса, а перенесённая
разметка не проходила полной независимой ручной верификации. Кроме того,
сопоставление с LLM-оценщиком выполнялось как опорное: Qwen2.5-72B применялась
к исходному англоязычному RAGBench, тогда как обученные классификаторы
оценивались на RAGBench-RU.

Дальнейшее развитие работы может включать ручную проверку части русскоязычной
разметки, запуск LLM-оценщика непосредственно на RAGBench-RU, расширение
оценки на исходно русскоязычные корпуса и развитие сервиса для прикладной
интеграции.

\subsection{Артефакты работы}

Открытые результаты работы включают:
\begin{itemize}
  \item набор данных: \url{https://huggingface.co/datasets/CMCenjoyer/ragbench-ru};
  \item опубликованные веса лучшей конфигурации оценщика (ModernBERT, длина~2048): \url{https://huggingface.co/CMCenjoyer/ru-trace-modernbert-2048};
  \item исходный код обучающего конвейера и анализа: директория проекта;
  \item Streamlit-сервис: \texttt{new\_streamlit.py};
  \item таблицы и графики экспериментального анализа: \texttt{analysis\_output/};
  \item публикацию тезисов доклада по теме оценки качества RAG-систем в русскоязычном
  домене на конференции <<Ломоносовские чтения~- 2026>>~\cite{moiseev2026ragru}.
\end{itemize}
